%==============================================================================
% INTRODUCTION
%==============================================================================
\chapter*{Introduction}
\addcontentsline{toc}{chapter}{Introduction}
\markboth{Introduction}{Introduction}

\section*{Context of the internship}

In a nuclear power plant, the management of the fuel rod governs the safety of
the installation. A fuel rod is a tube a few metres long, made of a stack of
uranium oxide pellets enclosed in a metallic cladding. It must operate for
several years under extreme conditions : temperatures above \SI{1000}{\kelvin},
intense irradiation, fuel swelling, cladding creep, and more. Predicting its
evolution is therefore essential in order to minimise risk.

\textbf{OFFBEAT} (\textit{OpenFOAM Fuel BEhavior Analysis Tool}) is one of the
computational codes dedicated to this field. Developed chiefly at the Laboratory
for Reactor Physics and Systems Behaviour of EPFL, at the Paul Scherrer
Institute and at Texas A\&M University, it is built on the open-source OpenFOAM
library and solves, using the finite volume method, the coupled thermal and
mechanical behaviour of a rod in one, two or three dimensions~[1]. Its use
follows the logic of OpenFOAM : a case is a tree of dictionary files that the
user must know how to write, a mesh to be generated, a solver to be launched
from the command line, and VTK-format outputs to be post-processed. The main
difficulty with OFFBEAT is that this chain of tasks, while powerful, is not very
accessible.

\section*{Subject and objective}

The objective set at the start of the internship was to \textit{``build an agent
on the nuclear fuel performance code OFFBEAT''}. The task was therefore to design
a system able to bridge a request expressed in natural language and the actual
execution of an OFFBEAT simulation, taking inspiration from the open-source
project \textbf{AutoFLUKA}~[3], which applies this idea to the particle
transport code FLUKA.

The project was carried out in two stages. The first consisted in building the
\textbf{base agent} : a supervisor built on a large language model, orchestrating
domain tools able to create a case, launch the solver, automatically repair
certain crashes and post-process the results. The second stage extended the
system towards a \textbf{digital twin} of the fuel rod. This meant monitoring
safety criteria as new data arrive, predicting threshold crossings, and
returning instantaneous predictions through a reduced-order model.

\section*{Outline of the report}

The report follows the progression of the work. \textbf{Chapter~1} presents the
scientific and technical context : the physics of the fuel rod, the OFFBEAT code,
and the principles of agents built on language models. \textbf{Chapter~2}
describes the architecture of the base agent and details its components, in
particular the self-healing loop, which is its most distinctive contribution.
\textbf{Chapter~3} sets out the extension towards the digital twin through the
five building blocks developed. \textbf{Chapter~4} gathers the results : a
reference simulation over a complete irradiation cycle, the validation of the
surrogate, and above all the systematic stress-testing of the agent's
self-diagnosis capabilities. \textbf{Chapter~5} returns to the difficulties
encountered and the false starts, whose analysis drove several important
corrections, and presents the internship schedule. The conclusion examines the
possible directions for continuing the work.

\clearpage

%==============================================================================
% CHAPTER 1
%==============================================================================
\chapter{Scientific and technical context}
\markboth{Chapter 1 : Scientific and technical context}{}

\section{The physics of the fuel rod}

\subsection{Anatomy and loadings}

A fuel rod in a pressurised water reactor (PWR) is a Zircaloy tube, called the
\textit{cladding}, roughly \SI{9.5}{\milli\metre} in outer diameter and a few
metres long, inside which a column of cylindrical uranium dioxide
($\mathrm{UO_2}$) pellets is stacked. Between pellet and cladding there
initially remains a radial gap of a few tens of micrometres, filled with
pressurised helium.

This gap is central to the behaviour of the rod. It constitutes a thermal
resistance that the heat must cross before reaching the cladding. A radial
temperature gradient of several hundred kelvins therefore appears between the
centre of the pellet and its periphery.

Under irradiation, several coupled phenomena modify this geometry :
\begin{itemize}
  \item the \textbf{densification} and then the \textbf{swelling} of the fuel,
    under the effect of solid and gaseous fission products;
  \item the \textbf{relocation} of the fragments of the cracked pellet;
  \item the \textbf{creep} of the cladding, which slowly deforms under the
    coolant pressure (about \SI{15.5}{\mega\pascal} in a PWR);
  \item the \textbf{release of fission gases} (xenon, krypton), which degrades
    the conductivity of the gap and raises the internal pressure.
\end{itemize}

The combined effect is a \textbf{progressive closure of the gap}. Once contact
is established, the pellet mechanically loads the cladding : this is the
\textbf{pellet-cladding mechanical interaction} (PCMI).

\subsection{Safety criteria}

The design of a fuel rod rests on quantified safety criteria.
Table~\ref{tab:criteres} summarises those relevant here, explicitly separating
the ones a fuel performance code computes from those belonging to other
disciplines. This distinction proved essential when scoping the digital twin
(§\ref{sec:cadrage}).

\begin{table}[htbp]
\centering
\caption[Safety criteria and coverage by OFFBEAT]{Fuel rod safety criteria and
their coverage by OFFBEAT. The values given are design orders of magnitude, to
be confirmed according to the rod and the burnup considered.}
\label{tab:criteres}
\small
\begin{tabularx}{\textwidth}{@{}>{\raggedright\arraybackslash}p{4.1cm}
                                >{\raggedright\arraybackslash}p{3.9cm}
                                X
                                c@{}}
\toprule
\textbf{Criterion} & \textbf{Indicative limit} & \textbf{Quantity} &
\textbf{OFFBEAT}\\
\midrule
Pellet centreline melting & $T < \SI{3113}{\kelvin}$ (fresh $\mathrm{UO_2}$;
  decreases with burnup) & Maximum temperature & yes\\
Cladding temperature (PCT, LOCA criterion) & $< \SI{1477}{\kelvin}$
  (\SI{1204}{\celsius}) & Temperature in the cladding & if transient modelled\\
Cladding strain (PCMI) & circumferential strain $< \SI{1}{\percent}$ &
  $\theta\theta$ strain & yes\\
Gap closure & margin $>0$ & Gap width & yes\\
Internal pressure (\textit{no-liftoff}) & $< P_{\text{coolant}}$
  ($\approx\SI{15.5}{\mega\pascal}$) & Plenum pressure & if FGR model enabled\\
Cladding oxidation (ECR) & $< \SI{17}{\percent}$ & Oxide thickness &
  model dependent\\
Departure from nucleate boiling (DNBR) & $> 1.3$ & --- & \textbf{no} (input)\\
\bottomrule
\end{tabularx}
\end{table}

Note that the DNBR belongs to core thermal-hydraulics, outside the domain of
OFFBEAT. For the code it is a boundary condition.

\section{The OFFBEAT code}

\subsection{Position among fuel performance codes}

Two distinct generations of codes coexist in fuel performance simulation.

The first gathers the so-called \textbf{1.5D} codes, which discretise the rod
into independent axial slices coupled only through the plenum and the internal
pressure. \textbf{FRAPCON} and its transient counterpart \textbf{FRAPTRAN}, and
\textbf{TRANSURANUS}, the European solution, belong to this category~[4, 5].
Their strength is speed and robustness : they are validated against very large
experimental databases and remain the reference tools for design studies. Their
limitation is geometric : azimuthal cracking, a pellet edge or a non-axisymmetric
gradient escape them by construction.

The second generation is \textbf{multidimensional}. \textbf{BISON},
\textbf{ALCYONE}, the French solution, and \textbf{OFFBEAT} belong to it. These
codes solve the rod in two or three dimensions and therefore capture local
effects, at the cost of a far higher computational burden (of the order of three
orders of magnitude compared with a 1.5D code).

OFFBEAT occupies a particular position in this landscape, and this is what
justifies its choice here : it is \textbf{multidimensional}, \textbf{open
source}, and built on OpenFOAM, a widely used library. Code-to-code comparisons
also show good agreement between OFFBEAT and BISON~[2].

\subsection{Solution principle}

On an unstructured mesh, OFFBEAT solves the coupling between :
\begin{itemize}
  \item a \textbf{thermal problem} : conduction in the pellet and the cladding,
    transfer across the gap by gas conduction and radiation;
  \item a \textbf{mechanical problem} : elasticity, plasticity and creep, in a
    cell-centred finite volume formulation;
  \item \textbf{behavioural models} : burnup, swelling, densification, fission
    gas release, pellet-cladding contact.
\end{itemize}

An OFFBEAT case inherits the structure of an OpenFOAM case : a
\texttt{0/} directory for the initial fields, \texttt{constant/} for the
properties and models, \texttt{system/} for the run control. To these is added a
\texttt{rodDict} file, specific to OFFBEAT, which describes the rod geometry
block by block and from which a script generates the mesh. The detail of this
tree is given in appendix~\ref{ann:cas}. Nothing in this chain is tolerant to
error, and whatever the source of a crash, the logs do not point to a physical
cause.

\section{Agents built on language models}

\subsection{Principle : supervisor and tools}

An \textit{agent}, in the sense used for large language models, is not a model
that answers questions but a model given the ability to \textbf{act}. The
mechanism, known as \textit{tool calling}, consists in describing to the model a
set of available functions (their name, purpose and arguments). Faced with a
request, the model then produces not text but a structured call : which tool to
invoke, with which parameters. The program executes the call and returns the
result, on which the model continues its reasoning. A sequence of such calls
then leads to an answer to the initial request.

This scheme has a major advantage in a scientific context : the language model
computes nothing itself. It invents neither a temperature nor a stress. It
decides which program to run and interprets what that program returns. The
physics remains entirely inside the solver.

\subsection{Related work and starting point}

The closest work is \textbf{AutoFLUKA}~[3], published at the end of \num{2024},
which automates the Monte Carlo simulation workflows of the particle transport
code FLUKA (a supervisor built on LangChain, domain tools for input file
generation, execution and post-processing, a self-healing loop in case of a
crash, and a documentary assistant based on semantic search). The motivations
behind the present work are the same as those of AutoFLUKA : the learning curve
of a scientific computing code is steep, and the workflows are time-consuming
and prone to human error.

AutoOFFBEAT therefore reuses the skeleton of AutoFLUKA, substituting the FLUKA
parts with OFFBEAT features. Thanks to this help with the implementation, the
effort could be concentrated on what is specific to nuclear fuel, namely the
physical content of the knowledge bases. Three differences are worth
highlighting, as they delimit the two pieces of work.

The \textbf{domain}. FLUKA is a transport code; its output quantities are count
rates and fluences. OFFBEAT is a thermomechanical code whose outputs are
continuous fields comparable to \textbf{quantified safety criteria}. This
difference makes possible a feature absent from AutoFLUKA : the automatic
analysis of safety margins, with its editable threshold base.

\textbf{Prediction}. AutoFLUKA remains reactive : it executes and post-processes.
The extension towards the digital twin developed here (chapter~3), in particular
the Gaussian-process surrogate, makes it possible to answer questions without
re-running the solver.

\textbf{Evaluation}. The literature on such agents essentially reports
successful case studies. A significant part of the present work consisted in
measuring the limits of the system, from the actual reach of self-healing to the
accuracy of tool selection (chapter~4).

\section{Scope adopted}

The scope was fixed as follows :
\begin{itemize}
  \item the base agent is the priority, and is built \textbf{one block at a
    time}, each tested before moving to the next;
  \item the code must remain \textbf{independent of the language model
    provider}, so that it can run entirely locally and at no cost;
  \item the digital twin strand is only started once the base agent is stable;
  \item the scale adopted is the \textbf{single fuel rod}.
\end{itemize}

\clearpage
